\documentclass[12pt]{article}
\usepackage[a4paper,margin=1in]{geometry}
\usepackage{setspace}

\title{AutoBib Extended References Test: Diverse Citation Types}
\author{AutoBib Test Suite}
\date{2025}

\begin{document}
\maketitle
\onehalfspacing

\section{Books / Textbooks}
1. Software Engineering: A Practitioner's Approach - McGraw-Hill:
https://www.mhprofessional.com/9780078022128/software-engineering-a-practitioners-approach

2. The LaTeX Companion - Addison-Wesley:
https://www.pearson.com/en-us/subject-catalog/p/the-latex-companion/P200000003293/9780201362992

3. Modern Information Retrieval: The Concepts and Technology behind Search - Pearson:
https://www.pearson.com/en-us/subject-catalog/p/modern-information-retrieval/P200000003621/9780136072240

\section{Conference Proceedings}
4. "Automatic Citation and Indexing of Scientific Literature" - ACM SIGIR:
https://dl.acm.org/doi/10.1145/275599.275643

5. "Scholarly Document Processing: A Graph-based Approach for Citation Analysis" - IEEE ICDE:
https://ieeexplore.ieee.org/document/9458421

6. "Understanding Software Testing Challenges" - ACM SIGSOFT:
https://dl.acm.org/doi/10.1145/3368089.3409695

7. "Empirical Studies of Refactoring" - ACM ICSE:
https://dl.acm.org/doi/10.1145/337802.337835

\section{Official Documentation}
7. LaTeX BibTeX Documentation - CTAN:
https://ctan.org/pkg/bibtex

8. Python Requests Library Documentation:
https://requests.readthedocs.io/en/latest/

9. NumPy Documentation - NumPy Developers:
https://numpy.org/doc/stable/

\section{Open-Source Repositories}
10. bibtexparser - Python BibTeX Parser Library:
https://github.com/sciunto-org/python-bibtexparser

11. scholarly - Academic data extraction from Google Scholar:
https://github.com/scholarly-python-package/scholarly

12. crossrefapi - Python client for Crossref API:
https://github.com/fabiobatalha/crossrefapi

\section{Technical Reports and Standards}

13. IEEE Citation Reference Guide - IEEE:
https://ieeexplore.ieee.org/document/7078186

14. RFC 8259 - The JavaScript Object Notation (JSON) Data Interchange Format:
https://www.rfc-editor.org/rfc/rfc8259.html

15. W3C HTML5 Specification:
https://www.w3.org/TR/html5/

\section{Software & GitHub (The Right Way)}
Direct GitHub links often fail in automation because they lack metadata. 
However, major projects publish in JOSS (Journal of Open Source Software) to get a DOI.

% 1. The "xarray" Python library (Hosted on GitHub, cited via JOSS DOI)
% EXPECT: @article (because JOSS is a journal), but it represents Software.
1. Xarray (Open Source Python Package):
https://doi.org/10.21105/joss.00002

% 2. Tidyverse (R Software)
2. Welcome to the Tidyverse:
https://doi.org/10.21105/joss.01469

% 3. Raw GitHub Link (Test Case)
% EXPECT: This will likely be SKIPPED or fail, because it has no DOI.
% This tests if your code gracefully handles non-academic URLs.
3. Linux Kernel (Raw GitHub):
https://github.com/torvalds/linux

\section{Conference Proceedings (Official Sources)}
Software Engineering research lives in conferences (ICSE, FSE, ASE).

4. On the naturalness of code (The original idea):
https://doi.org/10.1145/2337223.2337376

\section{Official Documentation & Standards}
Citing technical standards (like JSON or HTTP) is crucial for technical papers.

5. Hypertext Transfer Protocol Version 2 (HTTP/2):
https://www.rfc-editor.org/rfc/rfc7540.html




\end{document}
